% ID: SEL-D0073-OB-E03
% Tanda S073.
\textbf{Parte I.} Hay $150-90=60$ coches de gasolina, repartidos por igual. Entre los diésel, $90-72=18$ son usados:
\[
 \begin{array}{c|rr|r}
 &\text{Nuevos}&\text{Usados}&\text{Total}\\\hline
 \text{Diésel}&72&18&90\\
 \text{Gasolina}&30&30&60\\\hline
 \text{Total}&102&48&150
 \end{array}
\]
\textbf{a)} $P(\text{nuevo})=102/150=\boxed{17/25}$.

\textbf{b)} Al condicionar a que sea usado, el total de referencia es 48. Por tanto, $P(\text{diésel}\mid\text{usado})=18/48=\boxed{3/8}$.

\textbf{Parte II.} Para la media de una población normal con desviación típica conocida, el margen es $E=z_{1-\alpha/2}\sigma/\sqrt n$. Al $99\%$, el valor crítico es $z_{0{,}995}\approx2{,}5758$. Exigimos que el margen sea como máximo 2:
\[
 \frac{2{,}5758\cdot6}{\sqrt n}\le2
 \quad\Longrightarrow\quad n\ge(3\cdot2{,}5758)^2\approx59{,}71.
\]
Redondeamos hacia arriba para respetar el límite de error. El tamaño mínimo es $\boxed{60}$.
